% Options for packages loaded elsewhere
\PassOptionsToPackage{unicode}{hyperref}
\PassOptionsToPackage{hyphens}{url}
\documentclass[aspectratio=169,ignorenonframetext]{beamer}
\newif\ifbibliography
\usepackage{pgfpages}
\setbeamertemplate{caption}[numbered]
\setbeamertemplate{caption label separator}{: }
\setbeamercolor{caption name}{fg=normal text.fg}
\beamertemplatenavigationsymbolsempty
% remove section numbering
\setbeamertemplate{part page}{
  \centering
  \begin{beamercolorbox}[sep=16pt,center]{part title}
    \usebeamerfont{part title}\insertpart\par
  \end{beamercolorbox}
}
\setbeamertemplate{section page}{
  \centering
  \begin{beamercolorbox}[sep=12pt,center]{section title}
    \usebeamerfont{section title}\insertsection\par
  \end{beamercolorbox}
}
\setbeamertemplate{subsection page}{
  \centering
  \begin{beamercolorbox}[sep=8pt,center]{subsection title}
    \usebeamerfont{subsection title}\insertsubsection\par
  \end{beamercolorbox}
}
% Prevent slide breaks in the middle of a paragraph
\widowpenalties 1 10000
\raggedbottom
\AtBeginPart{
  \frame{\partpage}
}
\AtBeginSection{
  \ifbibliography
  \else
    \frame{\sectionpage}
  \fi
}
\AtBeginSubsection{
  \frame{\subsectionpage}
}
\usepackage{iftex}
\ifPDFTeX
  \usepackage[T1]{fontenc}
  \usepackage[utf8]{inputenc}
  \usepackage{textcomp} % provide euro and other symbols
\else % if luatex or xetex
  \usepackage{unicode-math} % this also loads fontspec
  \defaultfontfeatures{Scale=MatchLowercase}
  \defaultfontfeatures[\rmfamily]{Ligatures=TeX,Scale=1}
\fi
\usepackage{lmodern}
\ifPDFTeX\else
  % xetex/luatex font selection
\fi
% Use upquote if available, for straight quotes in verbatim environments
\IfFileExists{upquote.sty}{\usepackage{upquote}}{}
\IfFileExists{microtype.sty}{% use microtype if available
  \usepackage[]{microtype}
  \UseMicrotypeSet[protrusion]{basicmath} % disable protrusion for tt fonts
}{}
\makeatletter
\@ifundefined{KOMAClassName}{% if non-KOMA class
  \IfFileExists{parskip.sty}{%
    \usepackage{parskip}
  }{% else
    \setlength{\parindent}{0pt}
    \setlength{\parskip}{6pt plus 2pt minus 1pt}}
}{% if KOMA class
  \KOMAoptions{parskip=half}}
\makeatother
\usepackage{color}
\usetheme{Madrid}
\definecolor{TecBlue}{RGB}{0,51,102}
\setbeamercolor{structure}{fg=TecBlue}
\setbeamercolor{frametitle}{fg=TecBlue}
\setbeamercolor{title}{fg=white}
\setbeamertemplate{navigation symbols}{}
\usepackage{fancyvrb}
\newcommand{\VerbBar}{|}
\newcommand{\VERB}{\Verb[commandchars=\\\{\}]}
\DefineVerbatimEnvironment{Highlighting}{Verbatim}{commandchars=\\\{\}}
% Add ',fontsize=\small' for more characters per line
\newenvironment{Shaded}{}{}
\newcommand{\AlertTok}[1]{\textcolor[rgb]{1.00,0.00,0.00}{\textbf{#1}}}
\newcommand{\AnnotationTok}[1]{\textcolor[rgb]{0.38,0.63,0.69}{\textbf{\textit{#1}}}}
\newcommand{\AttributeTok}[1]{\textcolor[rgb]{0.49,0.56,0.16}{#1}}
\newcommand{\BaseNTok}[1]{\textcolor[rgb]{0.25,0.63,0.44}{#1}}
\newcommand{\BuiltInTok}[1]{\textcolor[rgb]{0.00,0.50,0.00}{#1}}
\newcommand{\CharTok}[1]{\textcolor[rgb]{0.25,0.44,0.63}{#1}}
\newcommand{\CommentTok}[1]{\textcolor[rgb]{0.38,0.63,0.69}{\textit{#1}}}
\newcommand{\CommentVarTok}[1]{\textcolor[rgb]{0.38,0.63,0.69}{\textbf{\textit{#1}}}}
\newcommand{\ConstantTok}[1]{\textcolor[rgb]{0.53,0.00,0.00}{#1}}
\newcommand{\ControlFlowTok}[1]{\textcolor[rgb]{0.00,0.44,0.13}{\textbf{#1}}}
\newcommand{\DataTypeTok}[1]{\textcolor[rgb]{0.56,0.13,0.00}{#1}}
\newcommand{\DecValTok}[1]{\textcolor[rgb]{0.25,0.63,0.44}{#1}}
\newcommand{\DocumentationTok}[1]{\textcolor[rgb]{0.73,0.13,0.13}{\textit{#1}}}
\newcommand{\ErrorTok}[1]{\textcolor[rgb]{1.00,0.00,0.00}{\textbf{#1}}}
\newcommand{\ExtensionTok}[1]{#1}
\newcommand{\FloatTok}[1]{\textcolor[rgb]{0.25,0.63,0.44}{#1}}
\newcommand{\FunctionTok}[1]{\textcolor[rgb]{0.02,0.16,0.49}{#1}}
\newcommand{\ImportTok}[1]{\textcolor[rgb]{0.00,0.50,0.00}{\textbf{#1}}}
\newcommand{\InformationTok}[1]{\textcolor[rgb]{0.38,0.63,0.69}{\textbf{\textit{#1}}}}
\newcommand{\KeywordTok}[1]{\textcolor[rgb]{0.00,0.44,0.13}{\textbf{#1}}}
\newcommand{\NormalTok}[1]{#1}
\newcommand{\OperatorTok}[1]{\textcolor[rgb]{0.40,0.40,0.40}{#1}}
\newcommand{\OtherTok}[1]{\textcolor[rgb]{0.00,0.44,0.13}{#1}}
\newcommand{\PreprocessorTok}[1]{\textcolor[rgb]{0.74,0.48,0.00}{#1}}
\newcommand{\RegionMarkerTok}[1]{#1}
\newcommand{\SpecialCharTok}[1]{\textcolor[rgb]{0.25,0.44,0.63}{#1}}
\newcommand{\SpecialStringTok}[1]{\textcolor[rgb]{0.73,0.40,0.53}{#1}}
\newcommand{\StringTok}[1]{\textcolor[rgb]{0.25,0.44,0.63}{#1}}
\newcommand{\VariableTok}[1]{\textcolor[rgb]{0.10,0.09,0.49}{#1}}
\newcommand{\VerbatimStringTok}[1]{\textcolor[rgb]{0.25,0.44,0.63}{#1}}
\newcommand{\WarningTok}[1]{\textcolor[rgb]{0.38,0.63,0.69}{\textbf{\textit{#1}}}}
\usepackage{longtable,booktabs,array}
\usepackage{caption}
\captionsetup[table]{skip=6pt}
\newcounter{none} % for unnumbered tables
\usepackage{calc} % for calculating minipage widths
\usepackage{caption}
% Make caption package work with longtable
\makeatletter
\def\fnum@table{\tablename~\thetable}
\makeatother
\setlength{\emergencystretch}{3em} % prevent overfull lines
\providecommand{\tightlist}{%
  \setlength{\itemsep}{0pt}\setlength{\parskip}{0pt}}
\usepackage{bookmark}
\IfFileExists{xurl.sty}{\usepackage{xurl}}{} % add URL line breaks if available
\urlstyle{same}
% fallback for those not using the hyperref driver hyperxmp:
\makeatletter
\@ifundefined{xmpquote}{\newcommand{\xmpquote}[1]{#1}}{}
\makeatother
\hypersetup{
  pdftitle={Regression Analysis},
  pdfauthor={Juliho Castillo Colmenares},
  hidelinks,
  pdfcreator={LaTeX via pandoc}}

\title{\texorpdfstring{Regression Analysis}{Regression Analysis}}
\author{\texorpdfstring{Juliho Castillo
Colmenares}{Juliho Castillo Colmenares}}
\subtitle{Linear, Multiple, and Diagnostic Modeling}
\institute{Tec de Monterrey}
\date{}

\begin{document}
\frame{\titlepage}

\begin{frame}{Today's Agenda}
\begin{enumerate}
\item
  Foundations of Simple Linear Regression (OLS)
\item
  ANOVA \& Intervals (Confidence vs. Prediction)
\item
  Residual Diagnostics \& Normality Assessment
\item
  Stepwise Variable Selection \& Information Criteria
\item
  Non-Linearity \& Polynomial Transformations
\item
  Multiple Linear Regression (Matrix Formulation \& VIF)
\item
  Aberrant Observations \& Heteroskedasticity (Cook's D, HC3)
\item
  Practical Case Study: Residential Property Valuation
\end{enumerate}
\end{frame}

\begin{frame}{Simple Linear Regression: The OLS Model}
\textbf{Population Equation:}
\[y_{i} = \beta_{0} + \beta_{1}x_{i} + \epsilon_{i},\quad i = 1,\ldots,n\]

where \(\mathbb{E}\left\lbrack \epsilon_{i} \right\rbrack = 0\),
\(\operatorname{Var}(\epsilon_{i}) = \sigma^{2}\), and
\(\operatorname{Cov}(\epsilon_{i},\epsilon_{j}) = 0\).

\textbf{Ordinary Least Squares Estimators:}
\[{\hat{\beta}}_{1} = \frac{\sum_{i = 1}^{n}\left( x_{i} - \overline{x} \right)\left( y_{i} - \overline{y} \right)}{\sum_{i = 1}^{n}\left( x_{i} - \overline{x} \right)^{2}} = \frac{S_{xy}}{S_{xx}},\quad\quad{\hat{\beta}}_{0} = \overline{y} - {\hat{\beta}}_{1}\overline{x}\]

\textbf{Gauss-Markov Theorem:} Under the classical assumptions, OLS
estimators are the \textbf{Best Linear Unbiased Estimators (BLUE)}.
\end{frame}

\begin{frame}{ANOVA \& Variation Decomposition}
\textbf{Total Variation Partitioning:}
\[\underset{\text{ SST}}{\underbrace{\sum\left( y_{i} - \overline{y} \right)^{2}}} = \underset{\text{ SSR (Explained)}}{\underbrace{\sum\left( {\hat{y}}_{i} - \overline{y} \right)^{2}}} + \underset{\text{ SSE (Residual)}}{\underbrace{\sum\left( y_{i} - {\hat{y}}_{i} \right)^{2}}}\]

\textbf{Goodness of Fit (\(R^{2}\)):}
\[R^{2} = \frac{\text{SSR}}{\text{ SST}} = 1 - \frac{\text{SSE}}{\text{ SST}}\]

\textbf{Overall Model \(F\)-Test:}
\[F = \frac{\text{MSR}}{\text{ MSE}} = \frac{\frac{\text{SSR}}{k}}{\frac{\text{ SSE}}{n - k - 1}} \sim F_{k,n - k - 1}\]
\end{frame}

\begin{frame}{Mean Response vs. Individual Prediction Intervals}
For a new observation at \(x_{0}\):

{\def\LTcaptype{none} % do not increment counter
\begin{longtable}[]{@{}
  >{\raggedright\arraybackslash}p{(\linewidth - 2\tabcolsep) * \real{0.5000}}
  >{\raggedright\arraybackslash}p{(\linewidth - 2\tabcolsep) * \real{0.5000}}@{}}
\toprule\noalign{}
\endhead
\begin{minipage}[t]{\linewidth}\raggedright
\textbf{1. Mean Response CI:}
\[{\hat{y}}_{0} \pm t_{\text{crit}} \cdot s\sqrt{\frac{1}{n} + \frac{\left( x_{0} - \overline{x} \right)^{2}}{S_{xx}}}\]

\begin{itemize}
\item
  Captures uncertainty in the \textbf{average response}
  \(\mathbb{E}\left\lbrack Y~|~X = x_{0} \right\rbrack\).
\item
  Narrows as sample size \(n \rightarrow \infty\).
\end{itemize}
\end{minipage} & \begin{minipage}[t]{\linewidth}\raggedright
\textbf{2. Individual Prediction PI:}
\[{\hat{y}}_{0} \pm t_{\text{crit}} \cdot s\sqrt{1 + \frac{1}{n} + \frac{\left( x_{0} - \overline{x} \right)^{2}}{S_{xx}}}\]

\begin{itemize}
\item
  Captures uncertainty in a \textbf{single new outcome}
  \(Y_{\text{new}}\).
\item
  Always strictly wider due to \(+ 1\) term (\(\sigma^{2}\)).
\end{itemize}
\end{minipage} \\
\bottomrule\noalign{}
\end{longtable}
}
\end{frame}

\begin{frame}{Multiple Linear Regression (Matrix Formulation)}
\textbf{Vector Model:}
\(\mathbf{y} = \mathbf{X}\mathbf{\beta} + \mathbf{\epsilon}\), with
\(\mathbf{\epsilon} \sim \mathcal{N}_{n\left( \mathbf{0},\sigma^{2}\mathbf{I}_{n} \right)}\).

\textbf{Normal Equations \& OLS Solution:}
\[\mathbf{X}^{T}\mathbf{X}\hat{\mathbf{\beta}} = \mathbf{X}^{T}\mathbf{y} \Longrightarrow \hat{\mathbf{\beta}} = \left( \mathbf{X}^{T}\mathbf{X} \right)^{- 1}\mathbf{X}^{T}\mathbf{y}\]

\textbf{Covariance of Estimates \& Hat Matrix:}
\[\operatorname{Cov}(\hat{\mathbf{\beta}}) = \sigma^{2}\left( \mathbf{X}^{T}\mathbf{X} \right)^{- 1},\quad\quad\hat{\mathbf{y}} = \mathbf{H}\mathbf{y} = {\mathbf{X}(\mathbf{X}^{T}\mathbf{X})}^{- 1}\mathbf{X}^{T}\mathbf{y}\]

The diagonal \(h_{ii}\) represents the \textbf{leverage} of observation
\(i\).
\end{frame}

\begin{frame}{Multicollinearity: The VIF Diagnostic}
\textbf{Problem:} High correlation between predictors inflates standard
errors \(\operatorname{Var}({\hat{\beta}}_{j})\).

\textbf{Variance Inflation Factor (VIF):}
\[\text{ VIF}_{j} = \frac{1}{1 - R_{j}^{2}}\] where \(R_{j}^{2}\) is the
\(R^{2}\) obtained from regressing \(X_{j}\) onto all other remaining
predictors.

\textbf{Interpretation Guidelines:}

\begin{itemize}
\item
  \(\text{VIF } = 1\): Completely independent features.
\item
  \(1 < \text{ VIF } < 5\): Acceptable mild correlation.
\item
  \(\text{VIF } \geq 10\): Severe collinearity; requires feature pruning
  or Ridge/PCA regularization.
\end{itemize}
\end{frame}

\begin{frame}{Aberrant Data: Outliers \& Cook's Distance}
{\def\LTcaptype{none} % do not increment counter
\begin{longtable}[]{@{}
  >{\raggedright\arraybackslash}p{(\linewidth - 2\tabcolsep) * \real{0.5000}}
  >{\raggedright\arraybackslash}p{(\linewidth - 2\tabcolsep) * \real{0.5000}}@{}}
\toprule\noalign{}
\endhead
\begin{minipage}[t]{\linewidth}\raggedright
\textbf{1. Studentized Residuals:}
\[r_{i} = \frac{e_{i}}{s\sqrt{1 - h_{ii}}}\]

\begin{itemize}
\item
  Identifies \(Y\)-space outliers (\(|r_{i}| > 3\)).
\end{itemize}

\textbf{2. Leverage (\(h_{ii}\)):}

\begin{itemize}
\item
  Identifies \(X\)-space extreme points
  (\(h_{ii} > \frac{2(k + 1)}{n}\)).
\end{itemize}
\end{minipage} & \begin{minipage}[t]{\linewidth}\raggedright
\textbf{3. Cook's Distance (\(D_{i}\)):}
\[D_{i} = \frac{r_{i}^{2}}{k + 1}\left( \frac{h_{ii}}{1 - h_{ii}} \right)\]

\begin{itemize}
\item
  Measures the overall influence of point \(i\) on \textbf{all} fitted
  values \(\hat{\mathbf{y}}\).
\item
  Threshold: \(D_{i} > \frac{4}{n}\) warrants detailed investigation.
\end{itemize}
\end{minipage} \\
\bottomrule\noalign{}
\end{longtable}
}
\end{frame}

\begin{frame}{Heteroskedasticity: Breusch-Pagan \& Robust SE}
\textbf{Violation:}
\(\operatorname{Var}(\epsilon_{i}~|~\mathbf{X}) = \sigma_{i}^{2} \neq \text{ constant}\).

\textbf{Breusch-Pagan Test:} Regress squared residuals
\(\frac{e_{i}^{2}}{{\hat{\sigma}}^{2}}\) on \(\mathbf{X}\). Test
statistic \(LM = \frac{1}{2}\text{SSR}_{\text{aux}} \sim \chi_{k}^{2}\).
Reject \(H_{0}\) if \(p < 0.05\).

\textbf{Remedies:}

\begin{enumerate}
\item
  \textbf{Heteroskedasticity-Consistent (HC3) Robust SE:}
  \[\operatorname{Cov}_{HC3}\left( \hat{\mathbf{\beta}} \right) = \left( \mathbf{X}^{T}\mathbf{X} \right)^{- 1}\left( \sum_{i = 1}^{n}\frac{e_{i}^{2}}{\left( 1 - h_{ii} \right)^{2}}\mathbf{x}_{i}\mathbf{x}_{i}^{T} \right)\left( \mathbf{X}^{T}\mathbf{X} \right)^{- 1}\]
\item
  \textbf{Weighted Least Squares (WLS):} Weight observations by
  \(w_{i} = \frac{1}{\sigma_{i}^{2}}\).
\end{enumerate}
\end{frame}

\begin{frame}[fragile]{Case Study: Housing Valuation Workflow}
\textbf{Scenario:} 1,000 residential transactions with 7 structural \&
geographic predictors.

\begin{Shaded}
\begin{Highlighting}[]
\ImportTok{import}\NormalTok{ statsmodels.formula.api }\ImportTok{as}\NormalTok{ smf}

\CommentTok{\# 1. Quadratic Polynomial Specification}
\NormalTok{formula }\OperatorTok{=} \StringTok{"sale\_price \textasciitilde{} sqft\_living + bedrooms + bathrooms + "} \OperatorTok{\textbackslash{}}
          \StringTok{"dist\_city\_center\_km + building\_grade + house\_age\_years + I(house\_age\_years**2)"}

\CommentTok{\# 2. Fit with Robust HC3 Covariance}
\NormalTok{model }\OperatorTok{=}\NormalTok{ smf.ols(formula, data}\OperatorTok{=}\NormalTok{df).fit()}
\NormalTok{robust\_model }\OperatorTok{=}\NormalTok{ model.get\_robustcov\_results(cov\_type}\OperatorTok{=}\StringTok{"HC3"}\NormalTok{)}
\BuiltInTok{print}\NormalTok{(robust\_model.summary())}
\end{Highlighting}
\end{Shaded}

\textbf{Findings:} Model explains \(88.4\%\) of price variance with
valid robust \(p\)-values.
\end{frame}

\begin{frame}{Summary \& Key Takeaways}
\begin{itemize}
\item
  \textbf{Linear Regression} is the foundational tool for parametric
  inference and prediction.
\item
  Always verify OLS assumptions: linearity, normality of residuals,
  homoskedasticity, and no extreme collinearity (\(\text{VIF } < 5\)).
\item
  When non-linear patterns exist, incorporate polynomial terms or
  logarithmic transforms.
\item
  Guard against influential outliers using Cook's distance.
\item
  When heteroskedasticity is present, employ \textbf{HC3 robust standard
  errors} to protect inferential validity.
\end{itemize}

{ \textbf{Next Chapter: Chapter 2 --- Multivariate Analysis Foundations}
}
\end{frame}

\end{document}




